%% kapa: exam=1 round=24 year=2013 period=2 form=A subject=accounting subject_ko=회계학 seq=20 number=20
%% source: Q-Net 감정평가사 기출문제 (cst00309) artlSeq=5073802 file=extracted/1008/5073802/2013년 제24회 감정평가사 1차시험 기출문제/2교시(A형).hwp
%% conversion: kordoc 4.24.51 -> tools/gen1.py (auto); review=proofread-sonnet; flags=none
\begin{question}{20}
건설업체 (주)감평은 20×1년 5월 1일 (주)대한과 도급계약을 체결하였다. (주)감평은 진행기준에 의해 수익과 비용을 인식하며 진행률은 발생한 누적계약원가를 추정총계약원가로 나눈 비율로 측정한다. 공사기간은 20×4년 12월 31일까지이다. 최초 계약금액은 ￦100,000이었으며, 계약금액의 변동내역, 원가 등에 관한 자료가 다음과 같을 때 20×3년 말 미청구공사 잔액은?
(단위: ￦) \par
{\small\centering
\begin{tabular}{|c|c|c|c|c|}
\hline
\makecell{연도} & \makecell{20×1} & \makecell{20×2} & \makecell{20×3} & \makecell{20×4} \\
\hline
\makecell{당기 계약금액의 증가분(공사변경, 보상금, 장려금)} & \makecell{0} & \makecell{0} & \makecell{20,000} & \makecell{10,000} \\
\hline
\makecell{누적발생\newline{}계약원가} & \makecell{20,000} & \makecell{45,000} & \makecell{68,000} & \makecell{86,000} \\
\hline
\makecell{각 연도말에 추정한 예상추가원가} & \makecell{60,000} & \makecell{45,000} & \makecell{17,000} & \makecell{-} \\
\hline
\makecell{대금청구액} & \makecell{30,000} & \makecell{40,000} & \makecell{20,000} & \makecell{40,000} \\
\hline
\makecell{대금회수액} & \makecell{20,000} & \makecell{30,000} & \makecell{20,000} & \makecell{60,000} \\
\hline
\end{tabular}\par}

\begin{choices}
  \item ￦3,000
  \item ￦4,000
  \item ￦5,000
  \item ￦6,000
  \item ￦7,000
\end{choices}
\end{question}
\answer{4}
